\documentclass[11pt]{article}
\usepackage[T1]{fontenc}
\usepackage{textcomp}
\usepackage[margin=0.75in]{geometry}
\usepackage{amsmath,amssymb,amsthm}
\usepackage{array,booktabs}
\usepackage{hyperref}
\setlength{\parindent}{0pt}
\newtheorem{theorem}{Theorem}
\theoremstyle{definition}
\newtheorem{definition}{Definition}
\title{Flow Proof Artifact: List-rev-cons-undec}
\author{Generated by \texttt{flow doc proof}}
\date{Flow Proof Book}
\begin{document}
\maketitle
Reversing an undecuple cons list.\\[0.5em]
\textbf{Source.} church — https://en.wikipedia.org/wiki/Reverse\_(list)\\[0.5em]
\bigskip
\noindent{\large \textbf{Derived fact 1.} \textit{rev of eleven cons ending in nil reverses element order} \hfill List \cdot reverse \cdot reverse of undecuple cons}\par
\smallskip\noindent\textit{Coordinate: List \cdot reverse \cdot reverse of undecuple cons}\par
\smallskip\noindent\textit{Source: church}\par
\smallskip\noindent\textit{Built on: reverse of cons is append of reverse tail, for reverse on List, reverse of decuple cons, for reverse on List}\par
\medskip
\noindent\textbf{Goal.} rev of eleven cons ending in nil reverses element order\par
\noindent$\displaystyle \forall a \in \mathbb{N} \forall b \in \mathbb{N} \forall c \in \mathbb{N} \forall d \in \mathbb{N} \forall e \in \mathbb{N} \forall f \in \mathbb{N} \forall g \in \mathbb{N} \forall h \in \mathbb{N} \forall i \in \mathbb{N} \forall j \in \mathbb{N} \forall k \in \mathbb{N}\quad rev(cons(a, cons(b, cons(c, cons(d, cons(e, cons(f, cons(g, cons(h, cons(i, cons(j, cons(k, nil)))))))))))) = cons(k, cons(j, cons(i, cons(h, cons(g, cons(f, cons(e, cons(d, cons(c, cons(b, cons(a, nil))))))))))))$\par
\medskip
\noindent
\begin{tabular}{@{} >{\raggedright\arraybackslash}p{0.06\textwidth} >{\raggedright\arraybackslash}p{0.47\textwidth} >{\raggedleft\arraybackslash}p{0.41\textwidth} @{}}
\textbf{\#} & \textbf{Proof} & \textbf{Mathematics} \\
\midrule
\textbf{1.} & We prove that reverse of undecuple cons for reverse on List. & \\[0.45em]
\textbf{2.} & We invoke the derived fact governing reverse on List: reverse of cons is append of reverse tail, for reverse on List (instantiated for a, cons(b, cons(c, cons(d, cons(e, cons(f, cons(g, cons(h, cons(i, cons(j, cons(k, nil))))))))))). & \\[0.45em]
\textbf{3.} & We invoke the derived fact governing reverse on List: reverse of decuple cons, for reverse on List (instantiated for b, c, d, e, f, g, h, i, j, k). & \\[0.45em]
\textbf{4.} & From step 2 and step 3, this implies rev(cons(a, cons(b, cons(c, cons(d, cons(e, cons(f, cons(g, cons(h, cons(i, cons(j, cons(k, nil)))))))))))) equals cons(k, cons(j, cons(i, cons(h, cons(g, cons(f, cons(e, cons(d, cons(c, cons(b, cons(a, nil)))))))))))). Hence proven. & $\displaystyle rev(cons(a, cons(b, cons(c, cons(d, cons(e, cons(f, cons(g, cons(h, cons(i, cons(j, cons(k, nil)))))))))))) = cons(k, cons(j, cons(i, cons(h, cons(g, cons(f, cons(e, cons(d, cons(c, cons(b, cons(a, nil))))))))))))$ \\[0.45em]
\bottomrule
\end{tabular}
\label{thm:List-reverse-reverse_of_undecuple_cons}
\par\medskip\hrule
\end{document}
